\chapter{Open Problems}
\label{ch:open}

Spivak's exercises are where his book happens. These are this book's. Each is an object in the ledger with status \emph{open} and, by the ledger's own rule, no evidence --- an open object that carried evidence of its claim would be rejected by the checker. Effort estimates come from the surveys in the research corpus and are as reliable as such estimates ever are.

\section{Days}

\begin{artifact}{OP.riscv.bitlogic256-drat}
Certify the 256-function table. The answer is known (\obj{M.riscv.bitlogic256}); only the proof is missing. For each function and each instruction set, a refutation of every chain of length one less than the minimum: on the order of a thousand variables per function, under a gigabyte in total.
\ArtifactScope{OP.riscv.bitlogic256-drat}
\end{artifact}

\begin{artifact}{OP.gpu.lop3-npn222}
Minimum \code{LOP3.LUT} count for all 222 four-input NPN classes. This is exact three-input-lookup-table synthesis on a shipping GPU instruction set, and nothing comparable has been published for it.
\ArtifactScope{OP.gpu.lop3-npn222}
\end{artifact}

\section{Weeks}

\begin{artifact}{OP.riscv.rgb565-lb7}
Prove or refute the specification's own lower bound: \emph{at least 7 instructions are needed to pack a 5:6:5 RGB value} on the base ISA. The highest rhetorical stake in this book.
\ArtifactScope{OP.riscv.rgb565-lb7}
\end{artifact}

\begin{artifact}{OP.riscv.zbb-table}
A certified table of minimum base-ISA sequence lengths for every Zbb, Zbkb and Zbs instruction. The ratified specification publishes prose motivations and not one instruction count. RISC-V's integer base has about forty instructions, no condition flags, and a formal specification; it is the cleanest instruction set on earth for this. Row by row.
\ArtifactScope{OP.riscv.zbb-table}
\end{artifact}

\begin{artifact}{OP.riscv.zicond-select3}
Is three the minimum for Zicond's conditional select? The specification says three with one temporary and never claims minimality; a live argument about whether the extension is worth its opcodes turns on it.
\ArtifactScope{OP.riscv.zicond-select3}
\end{artifact}

\begin{artifact}{OP.avx2.transpose4x4}
Is the eight-instruction four-by-four transpose minimal? It is the textbook SIMD shuffle --- in every game engine and every BLAS microkernel --- universally believed optimal, never proved. Sixteen tags across four registers, about 2,500 variables at depth eight.
\ArtifactScope{OP.avx2.transpose4x4}
\end{artifact}

\begin{artifact}{OP.avx2.perfectshuffle-certify}
Certify, or refute, the shuffle table LLVM ships for AArch64. It has about 6,561 entries, generated by a fixpoint over a hand-picked operator set with a non-additive cost model, and no certificate. Each entry is a few hundred variables; the whole table is embarrassingly parallel; and if even one entry is suboptimal, that is a headline.
\ArtifactScope{OP.avx2.perfectshuffle-certify}
\end{artifact}

\begin{artifact}{OP.wasm.i8x16-to-x64}
Minimal x64 lowerings for the canonical WebAssembly sixteen-byte shuffles. ARM64 does each in one instruction; x64 has no equivalent, and LLVM's canonicalization is known to break the browser engines' pattern matching. Three engines would consume the table.
\ArtifactScope{OP.wasm.i8x16-to-x64}
\end{artifact}

\begin{artifact}{OP.bilinear.rf2p6}
Is the bilinear complexity of six-term binary polynomial multiplication sixteen or seventeen? A 2026 lower bound of sixteen and a 2005 upper bound of seventeen; about seven thousand variables; symmetry, not size, is the wall. It underlies GHASH, which runs in every TLS stack.
\ArtifactScope{OP.bilinear.rf2p6}
\end{artifact}

\begin{artifact}{OP.sbox.keccak-chi5}
Is the bit-gate complexity of Keccak's five-bit S-box twelve or thirteen? The interval was published in 2023 alongside a remark that earlier optimality results in the same literature may be incorrect. A certified answer would be the first for any S-box.
\ArtifactScope{OP.sbox.keccak-chi5}
\end{artifact}

\begin{artifact}{OP.format.cost-model-artifact}
Ship the cost model as a checkable artifact. A claim of \emph{cost four under model M} is meaningless unless \(M\) is pinned and machine-readable, including the side conditions under which its additivity holds. Kilobytes; nearly free; nobody does it.
\ArtifactScope{OP.format.cost-model-artifact}
\end{artifact}

\section{Months}

\begin{artifact}{OP.aes.mixcolumns-lb}
A nontrivial lower bound on the XOR count of the AES MixColumns matrix. The best known count descended from ninety-four to eighty-eight over a decade; no lower bound beyond the trivial thirty-two has ever been published. A certified fifty-five would be the first.
\ArtifactScope{OP.aes.mixcolumns-lb}
\end{artifact}

\begin{artifact}{OP.unison.certify}
Certify one Unison schedule. Unison reports eighty-one percent of functions proven optimal under a constant-latency model and emits its constraint model in a standard format; a Rust constraint solver exists that emits optimality certificates checkable by a Rocq-verified checker, and it has never been pointed at a compiler.
\ArtifactScope{OP.unison.certify}
\end{artifact}

\begin{artifact}{OP.hd.p1-p25}
Checkable minimality certificates for the twenty-five Hacker's Delight synthesis benchmarks on a declared instruction set. The 2025 state of the art proves twenty-four of them cost-optimal with no artifact; an earlier system could not synthesize seven of them at all. Not one member has a checkable certificate on any instruction set.
\ArtifactScope{OP.hd.p1-p25}
\end{artifact}

\begin{artifact}{OP.format.closure-transcript}
A certificate format for universal reachability claims. \emph{Every four-bit permutation is reachable in at most fifteen gates} quantifies over \(16!\) objects; one refutation per instance is off by thirteen orders of magnitude. What is needed is a streaming closure transcript with a checkable frontier invariant and a cardinality check. Pilot it on \obj{M.riscv.table22}, whose answer is known.
\ArtifactScope{OP.format.closure-transcript}
\end{artifact}

\begin{artifact}{OP.kernel.word-prelude}
The problem underneath every other one in this book. Build the word of width \(n\) as a library on the zero-axiom naturals, so that theorems about words accumulate rather than being rebuilt per proof; then add proof by reflection --- a kernel-run enumerator for small cases and a verified refutation checker for large ones --- so that the theorems of Chapter~\ref{ch:machines} become kernel theorems rather than checked claims. Today they reconstruct as attestations: modules that assert the claim and its negation and contain none of the reasoning.
\ArtifactScope{OP.kernel.word-prelude}
\end{artifact}

\section{The catalogue}

The sixteen above are curated from roughly ninety in the research corpus that accompanies this book: open cases in finite combinatorics, in engineering mathematics from floating-point error bounds to cubature rules, in cryptographic gate complexity, in quantum circuit synthesis where the most-cited optimality claim in the field --- that the Toffoli gate needs seven T gates --- has no certificate either. Each survey ranks its candidates by the ratio of certifiability to effort and flags what it could not verify. The reader who wants a problem should start there; the reader who wants a method should start with the objects.
